% ==============================================================================
% PLANTILLA MAESTRA GENERAL - UCNL
% ==============================================================================
% Uso:
% 1. Duplica este archivo y renombra la copia por materia, semana y actividad.
% 2. Actualiza documenttitle, documentsubtitle, documentsubject, coursename,
%    coursecode, universitydepartment y datos de la figura docente.
% 3. Lee la planeacion semanal y NO pegues las instrucciones completas dentro del
%    producto final. Usa la planeacion como lista de cumplimiento.
% 4. Revisa la bibliografia indicada, toma notas, registra citas y redacta con
%    comprension propia.
% 5. Identifica el producto academico solicitado y construyelo dentro del
%    documento siempre que sea posible.
% 6. Entrega en PDF con la nomenclatura indicada en el aula.
%
% Formato institucional recomendado:
% - Fuente tipo Helvetica, equivalente visual a Arial.
% - Interlineado 1.5.
% - Citas autor-anio con natbib: usar \citep{clave} o \citet{clave}.
% - Referencias en bibliografia BibTeX, formato APA 7 en cuerpo y referencias.
% - Caratula, resumen breve, desarrollo, producto solicitado, conclusion y
%   bibliografia.
% - Redaccion academica clara, coherente, cohesionada, formal y sin faltas.
% - La portada puede llevar una marca de agua institucional discreta; ajustar
%   las variables coverwatermark... si se requiere apagarla o cambiar intensidad.
%
% Criterios generales de cumplimiento:
% - Responder exactamente a la actividad solicitada en la planeacion.
% - Identificar conceptos, elementos, criterios o categorias de forma precisa.
% - Analizar, interpretar y tomar postura; no solo copiar fuentes.
% - Organizar la informacion por jerarquia, secuencia y criterios claros.
% - Incluir conclusion personal cuando la actividad lo permita o lo solicite.
% - Usar citas y referencias en APA 7.
% - Evitar copiar las indicaciones de la actividad en el cuerpo final.
%
% Regla para productos visuales:
% - Si la actividad pide tabla, cuadro comparativo, esquema, mapa conceptual,
%   mapa mental, linea de tiempo, diagrama, matriz, lista de verificacion, red
%   conceptual, SQA, organizador grafico o producto similar, realizarlo
%   preferentemente con LaTeX/TikZ para mantener nitidez, evidencia editable y
%   diseno controlado.
% - Si la consigna exige Canva, Miro u otra herramienta externa, insertar una
%   captura nitida y conservar la explicacion dentro del documento.
% - Para tablas extensas, usar pagina limpia, sin encabezados ni pies, landscape,
%   letra pequena y restaurar el estilo al terminar.
%
% Criterios de redaccion personal:
% - No escribir para llenar espacio; escribir para sostener una idea.
% - Convertir informacion en criterio: que significa, como funciona, que problema
%   resuelve, que limite presenta y que consecuencia produce.
% - Estructura mental recomendada: problema -> sistema -> consecuencia -> postura.
% - Usar pensamiento de diagnostico: entrada -> proceso -> salida.
% - Cada parrafo debe tener funcion: presentar, explicar, comparar, valorar o
%   concluir.
% - Usar primera persona con moderacion academica: "Considero que",
%   "Desde esta perspectiva", "La posicion que sostengo es". Evitar "yo creo".
%
% Notas de enfoque personal segun enfoque_personal.txt:
% - La voz debe ser academica, tecnica, directa, estrategica y funcional.
% - El texto no debe sonar como estudiante pasivo que recopila: debe sonar como
%   alguien que interpreta, ordena, cruza datos y busca consecuencia practica.
% - Antes de redactar, identificar la funcion de cada seccion: abrir problema,
%   definir criterio, demostrar con evidencia, comparar, valorar o concluir.
% - Evitar frases genericas como "es importante porque si" o "desde siempre".
%   Sustituirlas por relaciones verificables: causa, efecto, limite, costo,
%   funcion, riesgo, implicacion o mejora.
% - Cada concepto debe pasar de definicion a analisis. Formula util:
%   concepto -> funcion -> riesgo si se aplica mal -> consecuencia -> postura.
% - La postura personal no debe ser opinion suelta. Debe tener tres piezas:
%   posicion, razon y efecto. Ejemplo: "Considero que X, porque Y; por ello Z".
% - Integrar experiencia o criterio propio sin perder formalidad: no narrar la
%   vida personal, sino usar mirada practica para explicar por que el tema
%   importa en educacion, derecho, instituciones, comunicacion o profesion.
%
% Estilo observado en actividades ya realizadas:
% - Abrir con una tesis concreta:
%   "El Derecho no solo ordena conductas: organiza poder, define limites y
%   distribuye consecuencias."
% - Formular el problema en negativo y positivo:
%   "El problema no es solo..., sino..."
% - Convertir el producto visual en argumento:
%   "El cuadro no solo recopila datos; reconstruye la funcion, el problema y el
%   punto que requiere correccion."
% - Cerrar con aprendizaje transferible:
%   "No basta con saber que dice la norma; tambien debe analizarse a quien
%   protege, a quien excluye y que razones permiten defenderla publicamente."
%
% Nivel cognitivo esperado:
% - Recordar: usar solo para definir conceptos basicos. No debe ser el nivel
%   dominante en una actividad ponderable.
% - Comprender: explicar con palabras propias y relacionar con el tema.
% - Aplicar: llevar el concepto a un caso, norma, texto, correo, dilema o hecho.
% - Analizar: separar elementos, detectar causas, efectos, tensiones y limites.
% - Evaluar: sostener una postura con criterios, fuentes y consecuencias.
% - Crear: elaborar producto propio: mapa, cuadro, linea de tiempo, ensayo,
%   propuesta, reformulacion, glosario, portafolio o solucion argumentada.
%
% Regla de nivel minimo:
% - Si la actividad vale puntos y pide producto, no quedarse en recordar o
%   comprender. Debe llegar por lo menos a aplicar y analizar; si pide conclusion
%   o postura, debe llegar a evaluar.
%
% Uso de las 100 tecnicas didacticas:
% - Esta plantilla conserva el manual "Tecnicas didacticas del aprendizaje" que
%   se incorporo en las plantillas de Filosofia del Derecho, Etica y Moral
%   juridica, y Redaccion en contextos virtuales.
% - Cuando la planeacion mencione una tecnica de la coleccion UnADM, identificar:
%   definicion de la tecnica, producto esperado, pasos de construccion, criterios
%   de revision y cierre reflexivo.
% - La tecnica didactica no debe verse como adorno. Debe mostrar seleccion,
%   jerarquia, relacion entre ideas y lectura propia.
% - Siempre que sea viable, construir la tecnica en LaTeX/TikZ: cuadro
%   comparativo, mapa conceptual, esquema, mapa mental, linea de tiempo, matriz,
%   red conceptual, diagrama de flujo, tabla SQA o lista de verificacion.
%
% Checklist antes de entregar:
% [ ] El titulo, materia, semana, docente y fecha estan actualizados.
% [ ] El objetivo coincide con la planeacion.
% [ ] El producto solicitado aparece visible, rotulado y completo.
% [ ] La introduccion plantea un problema o postura, no solo anuncia el tema.
% [ ] El desarrollo explica, conecta y analiza.
% [ ] Hay citas APA autor-anio dentro del texto.
% [ ] La bibliografia final coincide con las fuentes citadas.
% [ ] La conclusion responde que se comprendio y que postura se sostiene.
% [ ] El nivel cognitivo llega al menos a aplicacion y analisis.
% [ ] Si hay conclusion/postura, incluye posicion, razon y consecuencia.
% [ ] La tecnica didactica solicitada esta construida y explicada.
% [ ] No aparecen instrucciones copiadas de la planeacion dentro del producto.
% [ ] Ortografia, acentos, sintaxis y nombres propios revisados.
% [ ] Si se uso IA como apoyo, el texto fue comprendido, revisado y apropiado.
% ==============================================================================

% CREACION DEL DOCUMENTO
\documentclass[
	spanish,
	letterpaper, oneside
]{article}

% ==============================================================================
% INFORMACION DEL DOCUMENTO
% Actualizar en cada actividad.
% Ejemplos:
% \def\documenttitle {Cuadro comparativo de conceptos fundamentales}
% \def\documentsubtitle {Actividad 1 - Microeconomia}
% \def\documentsubject {Actividad 1}
% ==============================================================================
\def\documenttitle {Actividad 1 - Microeconomia}
\def\documentsubtitle {Actividad 1 - Microeconomia}
\def\documentsubject {Actividad 1}

\def\documentauthor {Martin Jonathan de la Cruz}
\def\coursename {Microeconomia}
\def\coursecode {MIC}

\def\universityname {Universidad Ciudadana de Nuevo Leon}
\def\universityfaculty {Programa academico UCNL}
\def\universitydepartment {Microeconomia}
\def\universitydepartmentimage {departamentos/logo-ucnl}
\def\universitydepartmentimagecfg {height=1.57cm}
\def\universitylocation {Monterrey, Nuevo Leon}

% INTEGRANTES, PROFESORES Y FECHAS
\def\authortable {
	\begin{tabular}{ll}
		\textbf{Alumno:}
		& \begin{tabular}[t]{l}
			 Martin Jonathan de la Cruz \\
		\end{tabular} \\
		Matricula: & UCNL \\

		\textit{Figura docente:}
		& \begin{tabular}[t]{l}
			Nombre \\ Apellido
		\end{tabular} \\
		& \\
		\multicolumn{2}{l}{Fecha de realizacion: \today} \\
		\multicolumn{2}{l}{\universitylocation}
	\end{tabular}
}

% IMPORTACION DEL TEMPLATE BASE
\input{template}

% FORMATO GENERAL
\usepackage{helvet}
\renewcommand{\familydefault}{\sfdefault}

% TABLAS Y PRODUCTOS VISUALES
\usepackage{array}
\usepackage{longtable}
\usepackage{pdflscape}
\usepackage{tikz}
\usetikzlibrary{
	arrows.meta,
	positioning,
	calc,
	matrix,
	fit,
	shapes.geometric,
	shadows.blur
}

% CITAS EN FORMATO AUTOR-ANIO, COMPATIBLE CON APA 7 EN EL CUERPO DEL TEXTO
% Usar:
% - Cita parentetica: \citep{clave}
% - Cita narrativa: \citet{clave}
% - Varias fuentes: \citep{clave1,clave2}
\setcitestyle{authoryear,open={(},close={)}}

% ==============================================================================
% AYUDAS DE REDACCION
% Quitar o comentar \pendiente cuando el documento final este terminado.
% ==============================================================================
\newcommand{\pendiente}[1]{\textcolor{red}{Refuerzo Ciclo A materializado}}

% ==============================================================================
% MARCA DE AGUA EN PORTADA
% - Se aplica solo a la portada porque usa \AddToShipoutPictureBG*.
% - Para apagarla en alguna actividad: \def\coverwatermarkenabled {false}
% - Usar opacidad baja para que no compita con titulo, datos ni logotipo.
% ==============================================================================
\def\coverwatermarkenabled {true}
\def\coverwatermarkimage {img/departamentos/logo-nuevo-leon.png}
\def\coverwatermarkopacity {0.16}
\def\coverwatermarkwidth {0.72\paperwidth}
\def\coverwatermarkxshift {0cm}
\def\coverwatermarkyshift {-0.35cm}

\newcommand{\insertcoverwatermark}{%
	\ifthenelse{\equal{\coverwatermarkenabled}{true}}{%
		\AddToShipoutPictureBG*{%
			\begin{tikzpicture}[remember picture,overlay]
				\node[
					opacity=\coverwatermarkopacity,
					inner sep=0pt,
					xshift=\coverwatermarkxshift,
					yshift=\coverwatermarkyshift
				] at (current page.center) {%
					\includegraphics[width=\coverwatermarkwidth]{\coverwatermarkimage}%
				};
			\end{tikzpicture}%
		}%
	}{}%
}

\begin{document}

% PORTADA
\insertcoverwatermark
\templatePortrait

% CONFIGURACION DE PAGINA Y ENCABEZADOS
\templatePagecfg
\onehalfspacing

% ==============================================================================
% RESUMEN / ABSTRACT
% Debe ser breve. Formula recomendada:
% Tema + objetivo + tecnica/producto + enfoque personal.
% ==============================================================================
\begin{abstractd}
	\textbf{Objetivo:} Identificar los conceptos introductorios de la microeconomia mediante la resolucion de un cuestionario de opcion multiple sobre eleccion del consumidor, produccion, elasticidad, estructuras de mercado y equilibrio competitivo.

	\textbf{Producto elaborado:} Cuestionario resuelto sobre principios basicos de microeconomia, con enfasis en oferta y demanda, costo de oportunidad, utilidad, competencia, monopolio e intervencion gubernamental.

	\textbf{Enfoque de analisis:} El cuestionario se resuelve desde una perspectiva funcional: cada reactivo se interpreta segun la decision economica que explica, el incentivo que pone en juego y la consecuencia que produce sobre eficiencia, bienestar y asignacion de recursos.

	\textbf{Nivel cognitivo:} Comprender y aplicar, porque los reactivos exigen reconocer conceptos clave y relacionarlos con situaciones economicas concretas de mercado, consumo y produccion.
\end{abstractd}

% TABLA DE CONTENIDOS - INDICE
\templateIndex

% CONFIGURACIONES FINALES
\templateFinalcfg

% ======================= INICIO DEL DOCUMENTO =======================

% ==============================================================================
% SECCION 1: INTRODUCCION
% Apertura recomendada:
% - Iniciar con una afirmacion clara, no con una frase generica.
% - Traducir la consigna a un problema academico, juridico, etico, comunicativo
%   o profesional, segun la materia.
% - Explicar por que importa.
% - Cerrar anticipando el enfoque personal.
% ==============================================================================
\section{Introduccion}

La microeconomia no estudia a la economia como un bloque uniforme, sino como una red de decisiones concretas tomadas por consumidores, empresas y agentes que enfrentan recursos escasos. Su valor academico radica en mostrar como esas decisiones, aunque parezcan individuales, producen resultados visibles en precios, cantidades, competencia y bienestar.

El problema que atiende esta actividad consiste en distinguir con precision conceptos que suelen confundirse cuando solo se memorizan definiciones: utilidad marginal, costo marginal, elasticidad, producto marginal, monopolio, competencia perfecta, dilema del prisionero o peso muerto. Comprender estos terminos permite interpretar por que una decision racional a nivel individual no siempre genera el mejor resultado colectivo y por que la estructura del mercado modifica los incentivos de productores y consumidores.

Desde esta perspectiva, el cuestionario se responde con un criterio de diagnostico: en cada reactivo se identifica primero que concepto economico esta en juego, despues que relacion explica y, por ultimo, que consecuencia practica produce sobre el equilibrio del mercado o el bienestar social.

% ==============================================================================
% SECCION 2: DESARROLLO / ANALISIS
% Adaptar los subtitulos segun la materia y la planeacion.
%
% Productos posibles:
% - Cuadro comparativo: criterios, evidencia, diferencia, valoracion.
% - Mapa conceptual: concepto central, conceptos secundarios, conectores claros.
% - Linea de tiempo: fecha/etapa, hecho, aporte, consecuencia.
% - Estudio de caso: situacion, elementos, problema, interpretacion, solucion.
% - Lista de verificacion: criterio observable, indicador, cumplimiento, mejora.
% - Glosario: concepto, definicion, funcion, ejemplo, fuente.
% - Reporte o ensayo: tesis, argumentos, evidencia, postura y cierre.
%
% Regla de parrafo:
% Entrada: concepto o problema.
% Proceso: explicacion, fuente y relacion.
% Salida: consecuencia, criterio o postura.
% ==============================================================================
\section{Desarrollo}

\subsection{Contexto y problema}

La Actividad 1 se ubica en la fase introductoria del curso, donde se espera que el estudiante reconozca el lenguaje basico de la microeconomia y lo diferencie del enfoque macroeconomico. El contexto del cuestionario obliga a pasar de la definicion aislada al uso correcto del concepto, porque cada pregunta vincula teoria economica con una decision concreta: consumir, producir, competir, intercambiar o regular.

El problema de analisis no es solo seleccionar una opcion correcta, sino entender por que las otras opciones no describen adecuadamente el fenomeno economico. Por ello, el cuestionario funciona como una verificacion conceptual de entradas, procesos y resultados: que decide el agente, con que restricciones lo hace y que efecto genera esa decision en el mercado.

\subsection{Marco conceptual}

Los conceptos centrales del cuestionario son escasez, costo de oportunidad, oferta, demanda, elasticidad, utilidad marginal, producto marginal, costo marginal y estructura de mercado. Todos ellos permiten analizar como se asignan los recursos cuando no es posible satisfacer todas las necesidades al mismo tiempo. La microeconomia se concentra en esa escala de decision y en los incentivos que orientan la conducta de consumidores y empresas.

Tambien son relevantes la restriccion presupuestaria, la tasa marginal de sustitucion, el excedente del consumidor y el peso muerto. Estos conceptos muestran que las elecciones no dependen solo del deseo de consumir, sino tambien del ingreso disponible, de los precios relativos y de la eficiencia con la que opera el mercado. Cuando la competencia funciona, tiende a mejorar el bienestar; cuando aparecen monopolios o intervenciones ineficientes, pueden surgir distorsiones, menor competencia y perdidas sociales.

\subsection{Nivel cognitivo aplicado}

El nivel cognitivo dominante es comprender y aplicar. Comprender implica reconocer que estudia la microeconomia, que representa una curva de posibilidades de produccion o que define a un monopolio; aplicar implica trasladar esos conceptos a situaciones de mercado para identificar elasticidad, escasez, equilibrio, competencia o efectos de una politica de precios.

En consecuencia, el producto solicitado no se limita a repetir definiciones. Cada reactivo exige clasificar el problema economico, interpretar su logica y seleccionar la opcion coherente con el principio teorico correspondiente.

\subsection{Producto solicitado}

\subsubsection{Cuestionario resuelto}

\begin{enumerate}
	\item Que estudia la microeconomia?\\
	Opciones: \textbf{a. Las decisiones individuales de consumidores y empresas} \quad b. El tipo de cambio \quad c. El crecimiento economico de largo plazo \quad d. El comportamiento agregado de la economia.

	\item Que representa una curva de posibilidades de produccion?\\
	Opciones: a. El crecimiento economico en 10 anos \quad b. El ingreso total de una empresa \quad c. La inflacion de un pais \quad \textbf{d. Las combinaciones de bienes que se pueden producir con los recursos disponibles}.

	\item Que es la diferenciacion de productos?\\
	Opciones: a. Reducir calidad \quad b. Aumentar impuestos \quad \textbf{c. Hacer que el producto sea distinto a los demas} \quad d. Eliminar competencia.

	\item Que es el producto marginal del trabajo?\\
	Opciones: a. Salario promedio \quad \textbf{b. Aumento en la produccion por un trabajador adicional} \quad c. Tiempo trabajado \quad d. Ingreso total.

	\item Que indica la tasa marginal de sustitucion?\\
	Opciones: a. El costo adicional de produccion \quad b. El salario por hora \quad c. El precio relativo entre dos bienes \quad \textbf{d. La preferencia por un bien sobre otro}.

	\item Que determina la ley de la oferta y la demanda?\\
	Opciones: a. La inversion extranjera \quad b. El tipo de cambio \quad c. Las tasas de interes \quad \textbf{d. El precio y cantidad de equilibrio en el mercado}.

	\item Cual de los siguientes es un tipo de estructura de mercado?\\
	Opciones: a. Importacion \quad b. Arancel \quad \textbf{c. Monopolio} \quad d. Deflacion.

	\item Que es una amenaza creible en un juego economico?\\
	Opciones: a. Una accion sin consecuencias \quad b. Una declaracion que nadie toma en serio \quad \textbf{c. Una estrategia que el jugador realmente ejecutaria si es necesario} \quad d. Una queja sin fundamento.

	\item Que ocurre en un dilema del prisionero clasico?\\
	Opciones: a. Todos obtienen su maxima ganancia \quad b. Los jugadores cooperan siempre \quad c. No hay conflicto \quad \textbf{d. La mejor decision individual lleva a un peor resultado conjunto}.

	\item Que es el producto marginal?\\
	Opciones: a. El ingreso por unidad vendida \quad b. El producto promedio \quad \textbf{c. El incremento del producto total por una unidad adicional de insumo} \quad d. El costo total dividido entre unidades.

	\item Que factores afectan la elasticidad de la demanda?\\
	Opciones: a. Nivel educativo \quad b. Solo el precio \quad \textbf{c. Tiempo, existencia de sustitutos, y tipo de bien} \quad d. Ingreso nacional.

	\item Que representa el peso muerto?\\
	Opciones: \textbf{a. Perdida de eficiencia en el mercado} \quad b. Ganancia no distribuida \quad c. Produccion extra \quad d. Precio de equilibrio.

	\item Que sucede con un bien elastico si sube su precio?\\
	Opciones: \textbf{a. La demanda baja mucho} \quad b. La demanda sube \quad c. No cambia \quad d. La oferta se reduce.

	\item Que relacion tiene la competencia con el bienestar economico?\\
	Opciones: a. No hay relacion \quad b. Lo reduce \quad \textbf{c. Generalmente lo incrementa} \quad d. Elimina el ingreso.

	\item Que es la utilidad marginal?\\
	Opciones: a. El gasto total en bienes \quad b. El ingreso extra por vender una unidad mas \quad \textbf{c. La satisfaccion adicional por consumir una unidad mas} \quad d. El costo de una unidad extra.

	\item Que representa el costo marginal?\\
	Opciones: a. El ingreso marginal \quad b. El gasto total \quad \textbf{c. El costo adicional por producir una unidad mas} \quad d. El costo fijo por unidad.

	\item Que mide la elasticidad cruzada de la demanda?\\
	Opciones: a. Cambios de precio \quad \textbf{b. Cambio en la demanda de un bien respecto al precio de otro bien} \quad c. Cambio en costos \quad d. Relacion entre dos mercados.

	\item Que sucede si los precios de los bienes bajan mientras el ingreso permanece igual?\\
	Opciones: \textbf{a. La linea de restriccion presupuestaria se desplaza hacia afuera} \quad b. Se reduce el consumo \quad c. La utilidad marginal desaparece \quad d. El consumidor entra en crisis.

	\item Que ocurre con el excedente del consumidor cuando el precio de mercado disminuye?\\
	Opciones: a. Permanece igual \quad b. Se transforma en excedente del productor \quad \textbf{c. Aumenta} \quad d. Disminuye.

	\item Que puede generar una intervencion gubernamental ineficiente?\\
	Opciones: a. Aumento de eficiencia \quad b. Reduccion de precios \quad \textbf{c. Menor competencia} \quad d. Bienestar total.

	\item En competencia perfecta, que relacion existe entre precio y costo marginal en el equilibrio?\\
	Opciones: a. No hay costo marginal \quad b. Precio menor que costo marginal \quad c. Precio mayor que costo marginal \quad \textbf{d. Precio igual a costo marginal}.

	\item Que define a un monopolio?\\
	Opciones: a. Muchas empresas vendiendo el mismo producto \quad b. Competencia libre \quad c. Muchos consumidores y vendedores \quad \textbf{d. Un solo vendedor que domina el mercado}.

	\item Que es un costo de oportunidad?\\
	Opciones: a. El dinero que se invierte en bolsa \quad \textbf{b. El valor de la mejor alternativa no elegida} \quad c. El beneficio obtenido por una empresa \quad d. El pago total por un bien.

	\item Que sucede en el largo plazo para una empresa?\\
	Opciones: a. No se pueden cambiar insumos \quad b. El producto marginal siempre aumenta \quad \textbf{c. Todos los factores son variables} \quad d. Todos los factores de produccion son fijos.

	\item Que efecto tiene un precio maximo impuesto por debajo del precio de equilibrio?\\
	Opciones: \textbf{a. Excedente de demanda (escasez)} \quad b. Excedente de oferta \quad c. Mercado en equilibrio sin efectos \quad d. Aumento de la eficiencia.

	\item Que determina la demanda de un bien?\\
	Opciones: a. Inversion extranjera \quad b. Solo el precio del bien \quad \textbf{c. Ingreso del consumidor, preferencias, precios relacionados} \quad d. Costos de produccion.

	\item Cual de los siguientes bienes tiene utilidad marginal decreciente?\\
	Opciones: \textbf{a. Cualquier bien en general} \quad b. Bienes gratuitos \quad c. Agua \quad d. Oro.

	\item Que es una restriccion presupuestaria?\\
	Opciones: a. El gasto en bienes publicos \quad \textbf{b. El limite de gasto segun ingreso y precios} \quad c. El gasto del gobierno \quad d. El ingreso que no se puede gastar.

	\item Cual es la funcion del empresario segun la microeconomia?\\
	Opciones: a. Reducir precios \quad \textbf{b. Coordinar los factores de produccion} \quad c. Supervisar inventarios \quad d. Administrar empleados.

	\item Que ocurre cuando se impone un precio maximo por debajo del equilibrio?\\
	Opciones: a. Excedente \quad \textbf{b. Escasez} \quad c. Superproduccion \quad d. Inflacion.
\end{enumerate}

\subsection{Tecnica didactica aplicada}

En esta actividad se utilizo la tecnica de cuestionario estructurado, porque permite verificar de manera puntual si el estudiante distingue conceptos, relaciones y consecuencias dentro del analisis microeconomico. La secuencia de reactivos hace visible la relacion entre decision individual, restriccion, intercambio y resultado de mercado, por lo que el producto no se limita a enlistar respuestas, sino que organiza el contenido para demostrar comprension operativa del tema.

% ==============================================================================
% BLOQUE OPCIONAL: TABLA O CUADRO COMPARATIVO AMPLIO
% Usar cuando el producto necesite mucho ancho:
% - \clearpage para iniciar en pagina limpia.
% - \pagestyle{empty} para quitar encabezado y pie en el producto amplio.
% - landscape para aprovechar el ancho horizontal.
% - scriptsize, tabcolsep y arraystretch para controlar legibilidad.
% - Restaurar \pagestyle{fancy} al terminar.
%
% Si la tabla es corta, puede omitirse landscape y usarse longtable normal.
% ==============================================================================
% \clearpage
% \pagestyle{empty}
% \begin{landscape}
% \begingroup
% \scriptsize
% \setlength{\tabcolsep}{4pt}
% \renewcommand{\arraystretch}{1.35}
% \begin{longtable}{@{}p{0.18\linewidth}p{0.39\linewidth}p{0.39\linewidth}@{}}
% \caption{Titulo del cuadro comparativo}\label{tab:cuadro-comparativo}\\
% \toprule
% \textbf{Criterio} & \textbf{Elemento 1} & \textbf{Elemento 2} \\
% \midrule
% \endfirsthead
% \toprule
% \textbf{Criterio} & \textbf{Elemento 1} & \textbf{Elemento 2} \\
% \midrule
% \endhead
% Concepto & \textit{Refuerzo Ciclo A: Contenido sintetico con cita si aplica.} & \textit{Refuerzo Ciclo A: Contenido sintetico con cita si aplica.} \\
% Caracteristicas & \textit{Refuerzo Ciclo A: Contenido.} & \textit{Refuerzo Ciclo A: Contenido.} \\
% Funcion & \textit{Refuerzo Ciclo A: Contenido.} & \textit{Refuerzo Ciclo A: Contenido.} \\
% Valoracion critica & \textit{Refuerzo Ciclo A: Contenido.} & \textit{Refuerzo Ciclo A: Contenido.} \\
% \bottomrule
% \end{longtable}
% \endgroup
% \end{landscape}
% \pagestyle{fancy}

% ==============================================================================
% BLOQUE OPCIONAL: PRODUCTO VISUAL EN TIKZ
% Usar para mapa conceptual, mapa mental, esquema, flujo, proceso o diagrama.
% Cada flecha debe expresar relacion; no usar nodos sueltos sin conectores.
% ==============================================================================
% \begin{figure}[H]
% 	\centering
% 	\begin{tikzpicture}[
% 		node distance=1.3cm,
% 		box/.style={rectangle, rounded corners=2pt, draw=black!65, fill=black!4, align=center, minimum width=3.2cm, minimum height=0.9cm},
% 		arrow/.style={-{Latex[length=2.5mm]}, thick, draw=black!65}
% 	]
% 		\node[box] (centro) {Concepto central};
% 		\node[box, right=of centro] (relacion) {Concepto relacionado};
% 		\node[box, below=of relacion] (efecto) {Consecuencia};
%
% 		\draw[arrow] (centro) -- node[above]{implica} (relacion);
% 		\draw[arrow] (relacion) -- node[right]{produce} (efecto);
% 	\end{tikzpicture}
% 	\caption{Titulo del producto visual.}
% \end{figure}

\subsection{Analisis propio}

El principal hallazgo del cuestionario es que muchos conceptos basicos de microeconomia se articulan alrededor de una misma pregunta: como asignan recursos los agentes cuando enfrentan restricciones. Esto se observa en reactivos sobre costo de oportunidad, restriccion presupuestaria, curva de posibilidades de produccion y utilidad marginal. La interpretacion que se desprende es que la microeconomia no describe solo precios, sino decisiones bajo escasez.

Tambien se advierte que la eficiencia del mercado depende de la estructura competitiva y de la calidad de los incentivos. Por ello aparecen temas como monopolio, amenaza creible, dilema del prisionero, peso muerto o precio maximo. La consecuencia de este enfoque es clara: comprender la teoria microeconomica ayuda a distinguir entre decisiones que mejoran el bienestar y decisiones que generan escasez, menor competencia o perdida de eficiencia.

% ==============================================================================
% SECCION 3: POSTURA PERSONAL
% Debe mostrar comprension real. No repetir fuentes; interpretar.
% ==============================================================================
\section{Postura personal}

Considero que la microeconomia es una herramienta de lectura practica de la realidad, porque permite entender que detras de cada precio, consumo, inversion o regulacion existe una estructura de incentivos que orienta la conducta de los agentes. La razon es que las decisiones economicas no ocurren en abstracto: dependen de costos, restricciones, sustitutos, competencia y expectativas. La consecuencia de comprender esto es que se pueden evaluar mejor tanto las decisiones empresariales como las politicas que afectan al mercado.

% ==============================================================================
% SECCION 4: CONCLUSION
% Cierre recomendado:
% - Responder directamente a la actividad.
% - Indicar que se comprendio, no solo que se elaboro.
% - Mencionar criterio personal final.
% - No introducir informacion nueva.
% ==============================================================================
% --- Ciclo A: refuerzo editorial materializado ---
\section*{Refuerzo editorial Ciclo A}
Este apartado consolida la mejora editorial incremental del nodo a partir del estándar maduro de AulaTeX. El refuerzo atiende brechas detectadas de bibliografía, citas, análisis propio, transferencia y cierre argumentado, sin sustituir la consigna original ni copiar contenido de los casos maestros.

\subsection*{Tesis de trabajo}
La actividad debe sostener una postura verificable: el producto solicitado no sólo organiza información, sino que permite interpretar el problema disciplinar, justificar decisiones y reconocer sus consecuencias académicas o profesionales.

\subsection*{Análisis propio}
El hallazgo central es que una entrega madura requiere pasar de la descripción a la interpretación. Por ello, el desarrollo debe explicar qué revela el producto construido, por qué es relevante para la materia y qué criterio permite evaluar su pertinencia. Esta operación convierte la evidencia reunida en argumento académico.

\subsection*{Transferencia profesional}
La transferencia consiste en vincular el producto con situaciones reales de desempeño: diagnóstico, toma de decisiones, comunicación de resultados, cumplimiento normativo, intervención educativa o resolución de problemas según el campo disciplinar. Así, la actividad adquiere valor formativo más allá de la plantilla.

\subsection*{Cierre argumentado}
En consecuencia, la calidad de la actividad depende de que el producto visible, las fuentes y la conclusión formen una secuencia coherente. La conclusión debe recuperar la tesis, indicar el principal aprendizaje y señalar una implicación práctica o conceptual.

% Brechas locales atendidas por Ciclo A: bibliografia_insuficiente, citas_insuficientes, pendientes_editoriales
% Reglas globales aplicadas:
% - Priorizar cierre de brecha recurrente: bibliografia_insuficiente (427 nodos).
% - Priorizar cierre de brecha recurrente: conclusion_argumentada (375 nodos).
% - Priorizar cierre de brecha recurrente: citas_insuficientes (328 nodos).
% - Priorizar cierre de brecha recurrente: analisis_propio (288 nodos).
% --- Fin Ciclo A ---

\section{Conclusion}

La Actividad 1 permitio resolver un cuestionario centrado en los conceptos fundamentales de la microeconomia y relacionarlos con decisiones reales de produccion, consumo e intercambio. A partir de los reactivos, se comprendio que la oferta y la demanda explican el equilibrio del mercado, que la utilidad y el costo marginal orientan decisiones racionales, y que la competencia o su ausencia modifican el bienestar economico. En sintesis, el cuestionario confirma que la microeconomia no solo define conceptos: explica como se toman decisiones y que efectos producen sobre la eficiencia y el bienestar social.

% ==============================================================================
% MANUAL DE TECNICAS DIDACTICAS
% Dejar incluido en la plantilla maestra. En una entrega final, puede comentarse
% esta linea si la docente solo solicita el producto de la semana.
% ==============================================================================
% Referencia opcional desactivada: archivo externo no disponible en este repositorio.

% BIBLIOGRAFIA
% Cambiar el nombre del archivo .bib segun la materia/actividad.
% Ejemplo: \bibliography{filosofia-del-derecho}
\bibliography{microeconomia}

% FIN DEL DOCUMENTO





































































































\subsection*{Citas de refuerzo Ciclo A}
La mejora editorial se vincula con la memoria documental disponible mediante las referencias \citep{unadm100TecnicasDidacticas2023}. Estas citas deben revisarse en una etapa disciplinar fina para sustituir o complementar fuentes genéricas por literatura específica del tema.

\end{document}


